% ============================================================================================
% J.-P. Benzécri, « Sur les variétés localement affines et localement projectives »,
% Bulletin de la S. M. F. 88 (1960), p. 229-332 (thèse Sc. math., Paris, 1960).
% DOI 10.24033/bsmf.1551 ; Numdam BSMF_1960__88__229_0.
%
% TRANSCRIPTION, not the author's source. Transcribed on 2026-11-20 by a Claude session
% (T-bz-chapter2, DEC-137) from the Numdam scan papers/benzecri/sources/bsmf-1960-88-229/benzecri1960.pdf
% (sha256 77e0b535382fa33109b7615fb9fa336151a86723ab8210b92b91c01b0e306306), reading the page
% images, with the scan's OCR text layer as a help only. TO BE CHECKED BY THE AUTHOR.
%
% Scope: Chapter 2 « Convexité de certaines variétés balayables » complete, journal pp. 246-259
% (PDF page = journal page - 227): from its title (p. 246, where Chapter 1 ends) to just before
% the title of Chapter 3 (p. 259).
%
% Conventions of the transcription (those of benzecri1960_ch1.tex).
%  * The French text is the printed text, misprints included; each misprint we noticed is kept and
%    flagged by a comment `% [sic] ...` giving the evident intended reading. Words cut at a line
%    end are rejoined. Nothing illegible was met; `[illisible]` would mark it.
%  * `% --- p. NNN ---` marks the start of journal page NNN.
%  * Each numbered statement is in its environment, the printed number as the optional argument
%    (with its prime, `1$'$`, and with the parenthesis printed after it, « (dans le cas affin) »).
%    Where a statement ends is not printed (the scan marks only where it begins): the end of each
%    environment is an editorial choice, made at the first paragraph that is commentary or
%    notation rather than part of the statement. Proofs (printed without a heading, usually
%    « En effet, ... », once « DÉMONSTRATION. — ») are in `proof` environments, idem.
%  * Unnumbered THÉORÈME, COROLLAIRE and REMARQUE blocks are `theoreme`, `corollaire` and
%    `remarque` environments without argument; the numbered REMARQUE 1 and 2 of § 1 are
%    `remarque` environments with the number as argument. Unnumbered EXEMPLE blocks are `exemple`
%    environments, which carry no label (they are not statements of the ledger).
%  * The list of properties 1°-7° of the pseudo-longueur (§ 2, p. 249), printed without a number
%    and cited throughout the chapter (« la septième propriété de la pseudo-longueur »,
%    « cf. 3°, 5° »), is an `enonce` environment: editorial, as the DEC-122 enonce of Chapter 1.
%    Its verification is printed after Déf. 3 and is interrupted by Prop. 2: two `proof`
%    environments, each preceded by a comment and carrying its target as an optional argument
%    `[Proof of \ref{...}]`, editorial and not printed (the only \ref outside a \bzcf).
%  * \bzcf{printed text}{\ref{...}} marks an explicit cross-reference the source makes to a
%    statement (« cf. chap. 1, § 1, déf. 14 », « la proposition 6 », « la septième propriété de la
%    pseudo-longueur »): it prints the printed text only; the \ref names the statement the PRINTED
%    text designates, even where that is a misprint (flagged [sic]). Chapter 1 statements are
%    named by their labels in legacy/labeled/benzecri1960_ch1.tex; Chapter 2 statements by the
%    labels of the scheme bz<kind>:c2s<section><initial><number>_<slug> listed in NOTES.md (the
%    \label's themselves are not in this file). A reference of a proof to its own statement is
%    not marked.
%  * $\R$ is the printed bold R (the reals); $\mathring{H}$ the printed H with a ring (an open
%    hemisphere); $\tilde{\varphi}$ the printed phi with a tilde; $M^0$, $N^0$ the printed M, N
%    with a raised circle; $)f, g($ the printed open interval.
% ============================================================================================
\documentclass[11pt]{article}
\usepackage[utf8]{inputenc}
\usepackage[T1]{fontenc}
\usepackage[french]{babel}
\usepackage{amsmath,amssymb}
\newcommand{\R}{\boldsymbol{R}}
\newcommand{\Rc}{\mathcal{R}}
\newcommand{\bzcf}[2]{#1}
% \bznum{n} prints " n" and nothing for an empty argument (unnumbered THÉORÈME, REMARQUE).
\newcommand{\bznum}[1]{\if\relax\detokenize{#1}\relax\else\ #1\fi}
\newenvironment{resume}[1][]{\par\medskip\noindent #1. }{\par\medskip}
\newenvironment{defin}[1][]{\par\medskip\noindent\textsc{Définition #1.} --- }{\par\medskip}
\newenvironment{prop}[1][]{\par\medskip\noindent\textsc{Proposition #1.} --- \itshape}{\par\medskip}
\newenvironment{theoreme}[1][]{\par\medskip\noindent\textsc{Théorème\bznum{#1}.} --- \itshape}{\par\medskip}
\newenvironment{lemme}[1][]{\par\medskip\noindent\textsc{Lemme #1.} --- \itshape}{\par\medskip}
\newenvironment{corollaire}{\par\medskip\noindent\textsc{Corollaire.} --- \itshape}{\par\medskip}
\newenvironment{remarque}[1][]{\par\medskip\noindent\textsc{Remarque\bznum{#1}.} --- }{\par\medskip}
\newenvironment{exemple}[1][Exemple]{\par\medskip\noindent\textsc{#1.} --- }{\par\medskip}
\newenvironment{enonce}{\par\medskip\noindent\itshape}{\par\medskip}
% The optional argument of `proof` is accepted and not printed (room for a ledger target).
\newenvironment{proof}[1][]{\par}{\par}

\begin{document}

% --- p. 246 ---
% (Chapter 1 ends at the top of p. 246; see benzecri1960_ch1.tex.)
\section*{Chapitre 2. Convexité de certaines variétés balayables.}

Voici le plan de ce chapitre :

On définit au paragraphe 1 la notion de poche pour les variétés localement projectives (resp.
affines).

Au paragraphe 2, on définit, sur l'ensemble des segments d'une variété localement projective, un
invariant projectif, positif ou nul, appelé pseudo-longueur.

Cet invariant est utilisé au paragraphe 3 pour démontrer le \bzcf{théorème
1}{\ref{bzthm:c2s3t1_poche_balayable}} sur les variétés balayables possédant une poche; deux
autres résultats, où la notion de poche n'intervient pas, sont ensuite démontrés (\bzcf{prop.
1}{\ref{bzprop:c2s3p1_projection_non_bornee}} et
\bzcf{2}{\ref{bzprop:c2s3p2_hemisphere_corps_convexe}}).
% « théorème 1 » as printed: the theorem of § 3 is itself printed without a number (« THÉORÈME. »).

\subsection*{1. Coques et poches.}

Nous distinguerons, pour plus de clarté, le cas affin du cas projectif.

\begin{defin}[1 (dans le cas affin)]
Soit $V$ une variété localement affine de dimension $n$ : on définit une \textit{poche} de $V$
par la donnée :

--- d'un ouvert connexe $P$ de $V$;

% --- p. 247 ---
--- d'une carte $\varphi$ de $P$ dans $\R^n$ compatible avec la structure affine de $P$;

--- d'une fonction linéaire $z$ sur $\R^n$ appelée cote.

Le triple $(P, \varphi, z)$ doit satisfaire aux deux conditions suivantes :

\textit{a}. Sur l'image $\varphi(P)$ la cote $z$ a pour borne inférieure $0$ et est majorée par
$1$;

\textit{b}. Soit $(M, N)$ un segment de $V$ dont l'extrémité $M$ appartienne à $P$. Soit
$\tilde{\varphi}$ une carte compatible d'un voisinage de $(M, N)$ dans $\R^n$ (\bzcf{\textit{cf.}
chap. 1, \S\,1, déf. 14}{\ref{bzdef:c1s1d14_segment_variete}}) ayant en $M$ même germe que
$\varphi$ (\bzcf{\textit{cf.} chap. 1, \S\,1, prop. 1,
2\textsuperscript{o}}{\ref{bzprop:c1s1p1_cartes_compatibles}}). Alors, si $z(\tilde{\varphi}(N))$
est inférieur à $1$, $N$ appartient à $P$.
\end{defin}

\begin{defin}[1$'$ (dans le cas projectif)]
Soit $V$ une variété localement projective de dimension $n$ : on définit une \textit{poche} de
$V$ par la donnée :

--- d'un ouvert connexe $P$ de $V$;

--- d'un hémisphère ouvert $\mathring{H}$ de $S(n)$;

--- d'une carte $\varphi$ de $P$ dans $S(n)$ compatible avec la structure projective, et telle
que
\[ \varphi(P)\subset\mathring{H}; \]

--- d'une fonction linéaire $z$, appelée cote, sur $\mathring{H}$ considéré comme espace affin
(\bzcf{\textit{cf.} chap. 1, \S\,1, déf. 6}{\ref{bzdef:c1s1d6_structure_affine_hemisphere}}).

(Le système $P, \mathring{H}, \varphi, z)$ doit satisfaire aux deux conditions suivantes :
% [sic] « (Le système P, H, φ, z) » : lire « Le système (P, H, φ, z) » (misplaced parenthesis).

\textit{a}. Sur l'image $\varphi(P)$, la cote $z$ a pour borne inférieure $0$ et est majorée par
$1$;

\textit{b}. Soit $(M, N)$ un segment de $V$ dont l'extrémité $M$ appartienne à $P$. Soit
$\tilde{\varphi}$ une carte compatible d'un voisinage de $(M, N)$ dans $S(n)$ ayant en $M$ même
germe que $\varphi$. Alors,
\[ \tilde{\varphi}(M, N)\subset\mathring{H} \]
et, si $z(\tilde{\varphi}(N))$ est inférieur à $1$, $N$ appartient à $P$.
\end{defin}

\begin{remarque}
Soient $V$ une variété localement affine de dimension $n$, $(P, \varphi, z)$ une poche de $V$, $j$
un isomorphisme compatible avec la structure affine de $\R^n$ sur un hémisphère ouvert
$\mathring{H}$ de $S(n)$ : le système $(P, \mathring{H}, j\circ\varphi, z\circ j^{-1})$ définit
une poche de $V$ considérée comme variété localement projective (\bzcf{\textit{cf.} chap. 1,
\S\,1, déf. 12, remarque}{\ref{bzrem:c1s1r1_affine_vers_projective}}).
\end{remarque}

\begin{exemple}[Exemple dans le cas affin]
Soient $W$ un ouvert étalé au-dessus de $\R^n$ et $p$ la projection de $W$ sur $\R^n$; soit $z$
une forme linéaire sur $\R^n$; soit $P$ une composante connexe de l'image réciproque par $p$ du
demi-espace ouvert : $z<1$.

Si, sur $p(P)$, $z$ a pour borne inférieure $0$, le triple $(P, p, z)$ définit une poche de $W$.

Il suffit de vérifier qu'est satisfaite la condition $b$. Soient $M\in P$ et $(M, N)$ un segment
de $W$; supposons $z(p(N))<1$. Démontrons que $N$ appartient
% --- p. 248 ---
à $P$; $p(M, N)$ est inclus dans le demi-espace ouvert $z<1$, donc $(M, N)$ est inclus dans une
composante connexe de l'image réciproque par $p$ de ce demi-espace ouvert, donc dans $P$,
composante connexe de $M$.
\end{exemple}

\begin{exemple}[Exemple dans le cas projectif]
Soient $W$ un ouvert étalé au-dessus de $S(n)$, $p$ la projection de $W$ sur $S(n)$,
$\mathring{H}$ un hémisphère ouvert de $S(n)$, $z$ une forme linéaire sur $\mathring{H}$ muni de
sa structure affine et $P$ une composante connexe dans $W$ de l'image réciproque par $p$ du
demi-espace ouvert $z<1$.

Si, sur $p(P)$, $z$ a pour borne inférieure $0$, et si, de plus, pour tout segment $(M, N)$ de
$W$ d'origine $M$ dans $P$, on a $p(MN)\subset\mathring{H}$, le système $(P, \mathring{H}, p, z)$
définit une poche de $W$. On vérifie, dans le cas affin, qu'est satisfaite la condition $b$.
\end{exemple}

Nous verrons (\bzcf{\S\,3, théorème}{\ref{bzthm:c2s3t1_poche_balayable}}) que, si $V$ est une
variété localement projective (resp. affine) connexe balayable, possédant une poche
$(P, \mathring{H}, \varphi, z)$ [resp. $(P, \varphi, z)$], on est toujours dans la situation
considérée en exemple. $\varphi$ se prolonge en une application projective (resp. affine) de $V$
sur $S(n)$ (resp. $\R^n$) et $P$ n'est autre qu'une composante connexe de l'image réciproque par
$\varphi$ du demi-espace ouvert de $\mathring{H}$ (resp. $\R^n$) d'équation $z<1$.

\begin{defin}[2]
Une poche $(P, \mathring{H}, \varphi, z)$ [resp. $(P, \varphi, z)$] est appelée \textit{coque} si
$\varphi(P)$ est un convexe.
\end{defin}

\begin{defin}[3]
Une coque $(P, \varphi, z)$ [resp. $(P, \mathring{H}, \varphi, z)$] d'une variété localement
affine (resp. projective) est appelée une \textit{coque stricte} si le convexe $\varphi(P)$ de
$\R^n$ [resp. si $\varphi(P)$ considéré comme convexe de l'espace affin $\mathring{H}$] ne contient
pas de droite.
\end{defin}

\begin{remarque}[1]
Il importe de distinguer entre droite de $\mathring{H}$, espace affin, et droite de $S(n)$ : une
droite de $\mathring{H}$ est une demi-droite de $S(n)$, ou portion de droite comprise entre deux
points diamétralement opposés.
\end{remarque}

\begin{remarque}[2]
Dire que la coque $(P, \mathring{H}, \varphi, z)$ est une coque stricte équivaut à dire que
$\overline{\varphi(P)}$ est un corps convexe de $S(n)$.
\end{remarque}

\subsection*{2. La pseudo-longueur.}

On définit sur l'ensemble des segments d'une variété localement projective un invariant projectif
appelé \textit{pseudo-longueur}. Nous avons vu (\bzcf{chap. 1, \S\,1, déf. 12,
remarque}{\ref{bzrem:c1s1r1_affine_vers_projective}}) que toute variété localement affine est munie
canoniquement d'une structure localement projective. La pseudo-longueur est ainsi définie pour les
variétés localement affines : c'est un invariant affin.

Nous allons énoncer les principales propriétés de la pseudo-longueur avant d'en donner la
définition.

% --- p. 249 ---
La pseudo-longueur d'un segment $(M, N)$ relativement à la variété localement projective $V$ est
notée :
\[ l_V(M, N). \]

% Wrapped in an enonce environment (editorial, as the DEC-122 enonce of Chapter 1): the list
% 1°-7° is printed without a number and is cited throughout the chapter.
\begin{enonce}
$l_V(M, N)$ possède les propriétés suivantes :

1\textsuperscript{o} \qquad $l_V(M, N)\geqq 0$; \qquad $l_V(M, M)=0$;

2\textsuperscript{o} Soit $g$ un automorphisme projectif de $V$, on a
\[ l_V(gM, gN)=l_V(M, N); \]

3\textsuperscript{o} Soit $(M, N)\subset(M', N')$, on a
\[ l_V(M, N)\leqq l_V(M', N'); \]

4\textsuperscript{o} Soit $V'$ un ouvert de $V$ considéré comme variété localement projective, on
a
\[ l_V(M, N)\leqq l_{V'}(M, N) \]
pour tout segment $(M, N)$ inclus dans $V'$;

5\textsuperscript{o} La pseudo-longueur est une fonction semi-continue supérieurement du segment,
c'est-à-dire que, si un segment $(M_i, N_i)$ tend vers $(M, N)$, on a
\[ \overline{\lim}\, l_V(M_i, N_i)\leqq l_V(M, N); \]

6\textsuperscript{o} Si $V$ est l'intérieur d'un corps convexe, $l_V(M, N)$ est une fonction
continue en $M$ et $N$ sur le produit $V\times V$; $l_V(M, N)$ est nulle si et seulement si
$M=N$;

7\textsuperscript{o} Soit $(P, \mathring{H}, \varphi, z)$ une poche de $V$; soit $M\in P$ un point
de $V$ tel que $z(\varphi(M))<\dfrac{1}{K}$ (où $K$ désigne un nombre réel supérieur à $1$); soit
$(M, N)$ un segment de $V$ dont la pseudo-longueur $l_V(M, N)$ soit inférieure à $\log K$; alors,
$N\in P$ et $z\,\varphi(N)<K\,z(\varphi(M))$.
% « z φ(N) » as printed (read: z(φ(N))).
\end{enonce}

La définition de la pseudo-longueur utilise la

\begin{defin}[1]
Soient $A$, $A'$, $B'$, $B$ quatre points d'un segment $(AB)$ de $S(n)$ écrits comme ils se
rencontrent si l'on ordonne de $A$ vers $B$ : on appelle \textit{mesure projective} de $A'B'$
\textit{relativement} à $AB$ le logarithme du rapport anharmonique $(A', B', B, A)$. Ce rapport
anharmonique est une quantité supérieure ou égale à l'unité; la notion de mesure projective est
invariante par les automorphismes de $GS(n, \R)$.
\end{defin}

\begin{prop}[1]
La mesure projective décroît si l'on remplace le segment $(A', B')$ par un segment contenu dans
$(A', B')$; elle décroît encore si l'on remplace le segment $AB$ par un segment le contenant.
\end{prop}
% (No proof of Prop. 1 is printed.)

% --- p. 250 ---
\begin{defin}[2]
Soit $(A', B')$ un segment d'une variété localement projective $V$; on dit que $(A', B')$ est
\textit{majoré par $k$ dans $V$} s'il existe un segment $(A, B)$ de $V$ contenant $(A', B')$ tel
que, si $\varphi$ est une carte compatible d'un voisinage de $(A, B)$ sur un ouvert de $S(n)$, on
ait : mesure projective de $\varphi(A', B')$ relativement à $\varphi(A, B)<k$.
\end{defin}

\begin{defin}[3]
La \textit{pseudo-longueur} $l_V(M, N)$ d'un segment $(M, N)$ d'une variété localement projective
est la borne inférieure des nombres $k$ qui majorent $(M, N)$ dans $V$.
\end{defin}

% Proof of the properties 1°-7° (the enonce of p. 249), first part; NOT a proof of Déf. 3.
% It is interrupted by Prop. 2 and resumes after the proof of Prop. 2. The optional argument
% (editorial, not printed) names the target for the ledger, which would otherwise attach this
% proof to the statement just above.
\begin{proof}[Proof of \ref{bzthm:c2s2_proprietes_pseudo_longueur}]
Nous allons vérifier que la pseudo-longueur ainsi définie possède les propriétés
1\textsuperscript{o} à 7\textsuperscript{o}.

Les propriétés 1\textsuperscript{o} et 2\textsuperscript{o} résultent trivialement de la
définition.

Les propriétés 3\textsuperscript{o} et 4\textsuperscript{o} résultent de \bzcf{la proposition
1}{\ref{bzprop:c2s2p1_mesure_projective_monotone}}.

La propriété 6\textsuperscript{o} résulte de la
\end{proof}

\begin{prop}[2]
Soient $\overline{V}$ un corps convexe de $S(n)$ d'intérieur $V$; $A$, $B$ deux points de $V$;
$F$, $F'$ les points d'intersection de la droite $AB$ avec la frontière de $\overline{V}$ :

$l_V(A, B)$ n'est autre que la mesure projective de $AB$ relativement à $FF'$;

$l_V(A, B)$ est une fonction continue sur $V\times V$ qui ne s'annule que sur la diagonale.
\end{prop}

\begin{proof}
La proposition résulte de ce que $F$, $F'$ dépendent continûment de $A$, $B$.
\end{proof}

% Proof of the properties 1°-7° (the enonce of p. 249), second part (5° and 7°); NOT a proof of
% Prop. 2 nor of Déf. 2. Optional argument editorial and not printed, as above.
\begin{proof}[Proof of \ref{bzthm:c2s2_proprietes_pseudo_longueur}]
Pour démontrer 5\textsuperscript{o}, nous démontrerons que, si $MN$ est un segment borné par $k$
dans $V$, il en est de même de tout segment assez voisin de $MN$. Soient $(A, B)$
(\bzcf{\textit{cf.} déf. 2}{\ref{bzdef:c2s2d2_segment_majore}}) un segment de $V$ et $\varphi$ une
carte compatible d'un voisinage $U$ de $(A, B)$ sur l'intérieur d'un corps convexe de $S(n)$ tels
que la mesure projective de $\varphi(M, N)$ relativement à $\varphi(A, B)$ soit inférieure à $k$.
Il suffit de montrer, d'après 4\textsuperscript{o}, que la pseudo-longueur $l_U$ de tout segment
suffisamment voisin de $MN$ est inférieure à $k$. Or, cela résulte aisément de \bzcf{la
proposition 2}{\ref{bzprop:c2s2p2_pseudo_longueur_corps_convexe}}.

Pour démontrer 7\textsuperscript{o}, nous calculerons en considérant la structure affine de
$\mathring{H}$.

Soient $F$, $F'$ les points extrémités d'un segment contenant $(A, B)$ tel que
\[ 1<(\varphi(A), \varphi(B), \varphi(F'), \varphi(F))<K, \]
où $\varphi$ désigne l'extension à un voisinage de $(F, F')$ de la carte $\varphi$ de $P$. Posons
\[ \varphi(A)=a, \qquad \varphi(B)=b, \qquad \varphi(F)=f, \qquad \varphi(F')=f'. \]
% As printed, the proof of 7° writes (A, B) for the segment (M, N) of 7°.
Considérons dans l'espace affin de $\mathring{H}$ le rapport de mesures algébriques de segments
\[ 1<\frac{af'\times bf}{af\times ff'}<K; \]
% [sic] « af × ff' » : lire « af × bf' » (the cross-ratio (a, b, f', f) = (af'/bf')(bf/af), as the
% two quotients af'/bf' and bf/af used just below show).
les quatre points $a$, $b$, $f$, $f'$ ont une cote $z$ positive.

% --- p. 251 ---
Si la cote de $b$ est inférieure ou égale à la cote de $a$, d'après la propriété $b$ de \bzcf{la
définition d'une poche}{\ref{bzdef:c2s1d1prime_poche_projective}}, $B\in P$, et la propriété
7\textsuperscript{o} est démontrée.

Si la cote de $b$ est supérieure à celle de $a$, on a les inégalités
\[ 0<z(f)<z(a)<z(b)<z(f'). \]
On en déduit que $\dfrac{af'}{bf'}$ est supérieur à l'unité, en sorte que le quotient
$\dfrac{bf}{af}$ est inférieur à $K$; d'où résulte l'inégalité
\[ z(b)<K\times z(a)<1 \]
et, d'après la propriété $b$ de \bzcf{la définition 1$'$ du paragraphe
1}{\ref{bzdef:c2s1d1prime_poche_projective}}, on a $B\in P$.
\end{proof}

\begin{defin}[4]
Soit $C$ une partie d'une variété localement projective $V$ : on dit que $C$ est majoré par $k$
dans $V$, si deux points quelconques de $C$ peuvent être joints par un segment qui est inclus dans
$C$ et est majoré par $k$ dans $V$.
\end{defin}

\begin{prop}[3]
Soit $V$ une variété localement projective, $k$ un nombre réel strictement positif. Tout point $M$
de $V$ possède un voisinage qui est majoré par $k$.
\end{prop}

\begin{proof}
Si $V$ est l'intérieur d'un corps convexe, la propriété est une conséquence facile de \bzcf{la
sixième propriété de la pseudo-longueur}{\ref{bzthm:c2s2_proprietes_pseudo_longueur}}.

Si $V$ n'est pas l'intérieur d'un corps convexe, choisissons un voisinage $U$ de $M$ dont il
existe une carte compatible dans l'intérieur d'un corps convexe. On peut, d'après ce qui précède,
trouver un voisinage $W$ de $M$, borné par $k$ dans $U$. D'après \bzcf{la quatrième propriété de
la pseudo-longueur}{\ref{bzthm:c2s2_proprietes_pseudo_longueur}}, $W$ est \textit{a fortiori}
majoré par $k$ dans $V$.
\end{proof}

\begin{defin}[5]
Soit $E$ une partie compacte d'une variété localement projective $V$; soit $k$ un nombre réel
strictement positif. On désigne par $n(E, k)$ le plus petit entier $n$, tel qu'il existe un
recouvrement de $E$, formé de $n$ ouverts majorés par $k$.
\end{defin}

\begin{prop}[4]
Soit $E$ une partie compacte connexe d'une variété localement projective $V$. Deux points $M$ et
$N$ quelconques de $E$ peuvent être joints par une ligne brisée formée d'au plus $n(E, k)$
segments majorés par $k$ dans $V$.
\end{prop}

\begin{proof}
Pour $n(E, k)=1$, la proposition 4 est équivalente à la définition : nous allons démontrer la
proposition par récurrence sur $n(E, k)$. Supposons la démonstration faite pour $n(E, k)<p$, et
supposons $n(E, k)=p$. Soient $O_1$, \ldots, $O_i$, \ldots, $O_p$, $p$ ouverts majorés par $k$
dont la réunion contienne $E$. Soit $M\in O_j$; $N\in O_k$; soit $O_Q=\bigcup_{i\in Q} O_i$, la
composante connexe contenant $N$, de la réunion des ouverts $O$ d'indice différent de $j$ : $O_Q$
est recouvert à
% --- p. 252 ---
l'aide d'au plus $p-1$ ouverts bornés par $k$; donc deux points quelconques de $O_Q$ peuvent être
joints par une ligne brisée formée d'au plus $p-1$ segments majorés par $k$ (hypothèse de
récurrence). D'autre part, puisque $O_a$ est composante connexe de la réunion des $O$ d'indice
différent de $j$, un tel ouvert est inclus dans $O_Q$, ou a avec $O_Q$ une intersection vide. $E$
étant connexe, il en résulte que $O_Q\cap O_j$ n'est pas vide; il suffira donc de joindre $M$ à un
point de cette intersection, puis ce point à $N$, pour relier $M$ à $N$ par une ligne brisée d'au
plus $p$ segments majorés par $k$.
% [sic] « O_a » : lire « O_Q ».
\end{proof}

\begin{defin}[6]
Soient $V$ une variété localement projective, $M\in V$, $r$ un nombre réel positif. On appelle
pseudo-boule de centre $M$ et de rayon $r$ l'ensemble des points $N$ de $V$ tels qu'il existe au
moins un segment fermé d'extrémités $M$, $N$, majoré par $r$.
\end{defin}

La pseudo-boule de centre $M$ et de rayon $r$, est un ouvert connexe (\bzcf{\textit{cf.}
3\textsuperscript{o}, 5\textsuperscript{o}}{\ref{bzthm:c2s2_proprietes_pseudo_longueur}}).

Si $V$ est l'intérieur d'un corps convexe, il résulte \bzcf{du
6\textsuperscript{o}}{\ref{bzthm:c2s2_proprietes_pseudo_longueur}} que lcs pseudo-boules de centre
$M$ et de rayon quelconque forment une base de voisinages de $M$.
% [sic] « lcs » : lire « les ».

\begin{prop}[5]
Soient $V$ une variété localement projective connexe, $U$ un ouvert non vide de $V$ et $r$ un
nombre réel, tels que, quel que soit $M\in U$, la pseudo-boule de centre $M$ et de rayon $r$ soit
incluse dans $U$. Alors $U=V$.
\end{prop}

\begin{proof}
Soit $M\in U$; $N\in V$ : nous allons montrer que $N\in U$. En effet, $M$ et $N$ peuvent être
joints par une ligne brisée formée de segments tous majorés par $r$ : on démontre de proche en
proche à partir de $M$ que tous les sommets de cette ligne appartiennent à $U$.
\end{proof}

\begin{prop}[6]
Soit $V$ une variété localement projective, possédant une poche $(P, \mathring{H}, \varphi, z)$.
Soient $A$ et $B$ deux points de $V$, $L$ une ligne brisée joignant $A$ à $B$, et formée de $p$
segments tous majorés dans $V$ par $k=\log K$. Supposons que $A\in P$; $z(\varphi(A))<1/K^p$ :
alors $B$ est aussi dans $P$, ainsi que $L$, et l'on a
\[ z(\varphi(B))<K^p\,z(\varphi(A)). \]
\end{prop}

\begin{proof}
La proposition se démontre par récurrence sur $p$ : pour $p=1$, c'est \bzcf{la septième propriété
de la pseudo-longueur}{\ref{bzthm:c2s2_proprietes_pseudo_longueur}}, le passage de $p$ à $p+1$ est
immédiat.
\end{proof}

\begin{prop}[7]
Soit $V$ une variété localement projective connexe et balayable pour un groupe $G$
d'automorphismes projectifs (\bzcf{\textit{cf.} chap. 1, \S\,1, déf.
15}{\ref{bzdef:c1s1d15_homogene_balayable_divisible}}). Soit $(P, \mathring{H}, \varphi, z)$ une
poche de $V$. Soit $E$ un compact de $V$. Il existe un automorphisme projectif $g\in G$, tel que
$gE\subset P$.
\end{prop}

\begin{proof}
\textsc{Démonstration.} --- Nous allons d'abord montrer que tout compact $E$ de $V$ est inclus
dans un domaine générateur connexe (pour le groupe $G$); nous
% --- p. 253 ---
démontrerons alors la proposition en supposant que $E$ est un domaine générateur compact connexe.
Soit $K$ un domaine générateur compact. $E\cup K$ est aussi un domaine générateur compact; puisque
$V$ est localement compact et localement connexe, on peut recouvrir $E\cup K$ au moyen d'un nombre
fini de compacts connexes; soit $E'$, la réunion de ces compacts : puisque $V$ est connexe, $E'$
est inclus dans un compact connexe $E''$; il suffit, pour obtenir $E''$, d'ajouter à $E'$ un
nombre fini de lignes brisées reliant les composantes connexes de $E'$.

Soit donc $E$ un domaine générateur compact connexe. Soit $k$ un nombre strictement positif
($k=\log K$). Soit $M^0$ un point de $P$ de cote $z\,\varphi(M^0)$ inférieure à $1/K^{n(E, k)}$.
Puisque $E$ est un domaine générateur pour $G$, il existe $g\in G$ automorphisme projectif de $V$
et $N^0\in E$, tels que $gN^0=M^0$. Soit $N$ un point quelconque de $E$ : nous allons montrer que
$gN\in P$.
% « Soit K un domaine générateur compact » and « k = log K » as printed: K names both.

En effet (\bzcf{prop. 4}{\ref{bzprop:c2s2p4_ligne_brisee}}), il existe une ligne brisée $L$
reliant $N^0$ à $N$, et formée d'au plus $n(E, k)$ segments tous bornés par $k$. Il résulte de
\bzcf{la proposition 6}{\ref{bzprop:c2s2p6_ligne_brisee_poche}} que $gN\in P$.
\end{proof}

\begin{corollaire}
Soit $V$ une variété localement projective, connexe balayable pour un groupe $G$
d'automorphismes projectifs; soit $(P, \mathring{H}, \varphi, z)$ une poche de $V$; soient $M$ et
$N$ deux points de $V$. Il existe un automorphisme projectif $g\in G$ de $V$ tel que $gM\in P$;
$gN\in P$.
\end{corollaire}

\begin{proof}
En effet, $\{M, N\}$ est une partie compacte de $V$.
\end{proof}

\begin{remarque}
Soit $V$ une variété localement affine balayable pour un groupe $G$ d'automorphismes affins.
Supposons que $V$ possède une poche $(P, \varphi, z)$. Pour la structure localement projective de
$V$ (\bzcf{\textit{cf.} chap. 1, \S\,1, déf. 12,
remarque}{\ref{bzrem:c1s1r1_affine_vers_projective}}), $G$ est un groupe d'automorphismes
projectifs; de plus (\bzcf{\textit{cf.} chap. 2, \S\,1, déf. 1$'$,
remarque}{\ref{bzrem:c2s1r0_poche_affine_vers_projective}}), si $j$ désigne un homéomorphisme de
$\R^n$ sur un hémisphère $\mathring{H}$ de $S(n)$ ($j$ compatible avec les structures affines),
$(P, \mathring{H}, j\circ\varphi, z\circ j^{-1})$ est une poche de la variété localement
projective de $V$. \bzcf{La proposition 7}{\ref{bzprop:c2s2p7_compact_dans_poche}} et \bzcf{son
corollaire}{\ref{bzcor:c2s2p7c_deux_points_dans_poche}} sont donc vrais aussi dans le cas affin :
on a l'énoncé :
\end{remarque}

\begin{prop}[7$'$]
Soit $V$ une variété localement affine, connexe et balayable pour un groupe $G$ d'automorphismes
affins; soit $(P, \varphi, z)$ une poche de $V$. Soit $E$ un compact de $V$. Il existe $g\in G$
tel que $gE\subset P$. Si $M$ et $N$ sont deux points de $V$, il existe $g\in G$ tel que :
$gM\in P$; $gN\in P$.
\end{prop}
% (No proof of Prop. 7' is printed: the REMARQUE just above is its justification.)

\subsection*{3. Applications.}

\begin{theoreme}
Soit $V$ une variété localement projective connexe balayable. Soit $(P, \mathring{H}, \varphi, z)$
une poche de $V$.

% --- p. 254 ---
1\textsuperscript{o} Il existe une carte compatible $\tilde{\varphi}$ de $V$ sur un ouvert de
$S(n)$, induisant $\varphi$ sur $P$. $P$ n'est autre que l'image réciproque par $\tilde{\varphi}$
de l'ensemble des points de $\mathring{H}$ où $z<1$.

2\textsuperscript{o} Si $P$ est simplement connexe, il en est de même de $V$.

3\textsuperscript{o} Si $\varphi(P)$ est convexe (\textit{i.~e.} si $(P, \mathring{H}, \varphi,
z)$ est une coque), $\tilde{\varphi}(V)$ est un convexe.

4\textsuperscript{o} Si $(P, \mathring{H}, \varphi, z)$ est une coque stricte,
$\tilde{\varphi}(V)$ est l'intérieur d'un corps convexe.
\end{theoreme}

\begin{proof}
La démonstration du 1\textsuperscript{o} va se décomposer en trois points.

\textit{a}. Il existe une application $\tilde{\varphi}$ de $V$ sur $S(n)$, coïncidant avec
$\varphi$ sur $P$, et dont le germe en tout point de $V$ appartient à $J(V)$ (\bzcf{\textit{cf.}
chap. 1, \S\,1, déf. 12}{\ref{bzdef:c1s1d12_variete_localement_projective}}).

Il faut montrer que le revêtement $J(V)$ de $V$ est isomorphe à un produit : $\tilde{\varphi}$ a
pour germes la composante connexe dans $J(V)$ des germes de $\varphi$.

Soit $[j, j']$ un chemin de $J(V)$, dont la projection sur $V$ par l'application source soit un
lacet $[M, M]$. Nous allons montrer que $j=j'$. D'après \bzcf{la proposition 7 du paragraphe
2}{\ref{bzprop:c2s2p7_compact_dans_poche}}, il existe un automorphisme projectif $g$ de $V$, tel
que l'image par $g$ du lacet $[M, M]$ soit dans $P$,
\[ g[M, M]\subset P. \]
Considérons maintenant le lacet de $J(V)$ obtenu à partir du lacet $[j, j']$ et composant chaque
jet avec l'automorphisme $g^{-1}$ : on obtient un lacet de $J(V)$ se projetant par l'application
source sur le lacet $g[M, M]$, inclus dans $P$. Mais $P$ est isomorphe par $\varphi$ à un ouvert
de $S(n)$ : donc la restriction à $P$ de $J(V)$ est un produit, et tout chemin de $J(V)$ se
projetant sur le lacet $g[M, M]\subset P$ est fermé. Le chemin $[j, j']$ dont l'image par
l'automorphisme $g^{-1}$ de $J(V)$ est fermé est donc fermé. L'assertation $a$ en résulte
immédiatement.
% [sic] « assertation » : lire « assertion ».

\textit{b}. L'application définie en $a$ de $V$ sur $S(n)$ est une carte compatible.

$\tilde{\varphi}$ est un homéomorphisme local : il faut montrer que $\tilde{\varphi}$ est un
homéomorphisme global, c'est-à-dire que deux points distincts $M$ et $N$ de $V$ ont par
$\tilde{\varphi}$ des images distinctes.

Soit $g$ un automorphisme projectif de $V$ : nous désignerons par $\gamma$ l'automorphisme de
$S(n)$ ($\gamma\in GS(n, \R)$) tel que
\[ \tilde{\varphi}\circ g=\gamma\circ\tilde{\varphi} \]
[l'ensemble des jets de $\tilde{\varphi}\circ g$ est une section de $J(V)$ au-dessus de $V$ :
puisque $GS(n, \R)$ opère transitivement sur la fibre des jets de même source, il existe $\gamma$
transformant $\tilde{\varphi}$ en la section $\tilde{\varphi}\circ g$].

Choisissons $g$ tel que $gM\in P$; $gN\in P$. Puisque $M\neq N$, on a $gM\neq gN$; et, puisque la
restriction de $\tilde{\varphi}$ à $P$ est un homéomorphisme, on a
\[ \tilde{\varphi}\circ gM\neq\tilde{\varphi}\circ gN. \]
% --- p. 255 ---
Mais il en résulte que
\[ \gamma\circ\tilde{\varphi}M\neq\gamma\circ\tilde{\varphi}N \]
et donc que
\[ \tilde{\varphi}M\neq\tilde{\varphi}N. \]
ce que nous voulions démontrer.

\textit{c}. $P$ est l'image réciproque par $\tilde{\varphi}$ de l'ensemble $D$ des points de
$\mathring{H}$, où $z<1$.

Montrons que
\[ P=\tilde{\varphi}^{-1}(D\cap\tilde{\varphi}(V)). \]
Soit $M\in\tilde{\varphi}^{-1}(D\cap\tilde{\varphi}(V))$ nous : allons montrer que $M\in P$. Soit
$N\in P$; on peut construire dans $D\cap\tilde{\varphi}(V)$ une ligne brisée allant de
$\tilde{\varphi}(M)$ à $\tilde{\varphi}(N)$. De la condition $b$ de \bzcf{la définition des
poches}{\ref{bzdef:c2s1d1prime_poche_projective}}, il résulte que $M$ appartient à $P$.
% [sic] « nous : allons montrer » : lire « nous allons montrer » (misplaced colon).

Dans la suite, nous supposerons, pour simplifier les notations, que $\tilde{\varphi}$ est
l'identité, et donc que $P$ et $V$ sont des parties de $S(n)$.

\textit{Démontrons} 2\textsuperscript{o}. --- Soit $L$ un lacet de $V$ : il faut montrer que $L$
est homotope à un point. Or il existe un automorphisme de $V$ qui envoie le compact $L$ dans $P$ :
comme $P$ est simplement connexe, $L$ est homotope à un point.

\textit{Démontrons} 3\textsuperscript{o}. --- Soit $M$ et $N$ deux points de $V$ : il faut
montrer que, dans $V$, il existe un segment fermé d'extrémités $M$, $N$. Or, soit $g$ un
automorphisme projectif tel que $gM\in P$, $gN\in P$ : il existe un segment $(gM, gN)$ puisque
$\varphi(P)$ est convexe; donc il existe un segment $(M, N)$.

\textit{Démontrons} 4\textsuperscript{o}. --- Si $V$ n'est pas l'intérieur d'un corps convexe, on
peut trouver dans $V$ une demi-droite $L$, c'est-à-dire une portion de droite limitée à deux
points diamétralement opposés de la frontière de $V$.

Soit $g$ un automorphisme projectif de $V$ qui envoie dans $P$ un point $M$ de $L$ : il résulte de
la condition $b$ de \bzcf{la définition des poches}{\ref{bzdef:c2s1d1prime_poche_projective}} que
$gL$ est une demi-droite incluse dans $\mathring{H}$, c'est-à-dire une droite de la géométrie
affine de $\mathring{H}$ : nous allons montrer que cette inclusion est incompatible avec
l'hypothèse faite que $P$ ne contient pas de droite de $\mathring{H}$.

En effet (\bzcf{\S\,1, déf. 1$'$, \textit{b}}{\ref{bzdef:c2s1d1prime_poche_projective}}), tout
point de $gL$ de cote $z$ inférieure à $1$ appartient à $P$ : donc tous les points de $gL$ ont une
cote supérieure ou égale à $0$, puisque la borne inférieure de $z$ sur $P$ est $0$. Mais alors,
$gL$ est une droite (pour la géométrie affine de $\mathring{H}$) de cote constante égale à $z(M)$
et $gL$ est incluse dans $P$, ce qui contredit l'hypothèse faite sur $P$.
% [sic] « z(M) » : lire « z(gM) » (M is the point of L sent into P by g).

La démonstration du théorème est achevée.
\end{proof}

Dans le cas affin, on a le

\begin{theoreme}
Soit $V$ une variété localement affine connexe balayable. Soit $(P, \varphi, z)$ une poche de $V$.

% --- p. 256 ---
1\textsuperscript{o} Il existe une carte compatible $\tilde{\varphi}$ de $V$ sur un ouvert de
$\R^n$. $P$ n'est autre que l'image réciproque par $\tilde{\varphi}$ de l'ensemble des points où
$z<1$.

2\textsuperscript{o} Si $P$ est simplement connexe, il en est de même de $V$.

3\textsuperscript{o} Si $(P, \varphi, z)$ est une coque, $\tilde{\varphi}(V)$ est un convexe.

4\textsuperscript{o} Si $(P, \varphi, z)$ est une coque stricte, $\tilde{\varphi}(V)$ ne contient
pas de droite.
\end{theoreme}

\begin{proof}
Le théorème se démontre \bzcf{comme le précédent}{\ref{bzthm:c2s3t1_poche_balayable}} : il suffit
de remplacer projectif par affin, et $S(n)$ ou $\mathring{H}$ par $\R^n$.
\end{proof}

\begin{prop}[1]
Soit $V$ une variété localement affine balayable munie d'une projection $p$ sur $\R^n$ compatible
avec la structure affine ($V$ est un ouvert étalé au-dessus de $\R^n$). Soit $E$ un domaine
générateur compact de $V$; soit $K$ un compact de $\R^n$. Il existe un segment $(M, N)$ de $V$,
tel que
\[ M\in E, \qquad p(N)\notin K. \]
\end{prop}

\begin{proof}
\textit{a}. La proposition est évidente si $V$ contient un fermé $L$ homéomorphe par $p$ à une
droite de $\R^n$.

En effet, il existe un automorphisme affin $g$ de $V$ tel que $gL\cap E\neq\emptyset$. On a
évidemment $p\circ gL\cap\complement K\neq\emptyset$. Il suffit de prendre
\[ M\in gL, \qquad p(N)\in p\circ gL\cap\complement K. \]

\textit{b}. Supposons maintenant que $V$ ne contienne pas de fermé $L$ homéomorphe par $p$ à une
droite de $\R^n$.

Nous décomposerons la démonstration de la proposition en une suite de points.

1\textsuperscript{o} Soit $k$ un nombre réel positif. Il existe un segment $(A, B)$ de $V$
jouissant de la propriété suivante :

Soit $(M, B)$ un segment de $V$ tel que
\[ (A, B)\subset(M, B). \]
Alors on a
\[ \frac{\overline{ma}}{\overline{ab}}<k, \]
où $m=p(M)$; $a=p(A)$; $b=p(B)$.

\textit{Démonstration} du 1\textsuperscript{o}. --- Soit $l$ une droite de $\R^n$ telle que
$p(V)\cap l\neq\emptyset$. Une composante connexe $S$ de $p^{-1}(l)$ se projette suivant un ouvert
connexe de $l$ qui pourra être un segment ouvert $)f, g($ ou une demi-droite $)f, \infty($.

Soit $(a, b)$ un segment fermé de $l$ tel que
\[ \frac{\overline{fa}}{\overline{fb}}<k. \]
% --- p. 257 ---
L'image réciproque par $p$ dans $S$ du segment $(a, b)$ est un segment $(A, B)$ de $V$ qui possède
les propriétés demandées.

2\textsuperscript{o} Supposons désormais $\R^n$ muni d'une métrique euclidienne. Il existe un
nombre $r$ positif tel que tout point de $E$ admette un voisinage $U$, tel que la restriction de
$p$ à $U$ définisse une carte sur une boule fermée de rayon $r$.

L'assertion résulte de la compacité de $E$.

3\textsuperscript{o} Il existe un segment $(A', B')$ de $V$ dont l'origine $A'$ appartient à $E$
et qui se projette par $p$ sur un segment de longueur supérieure à $r/k$.

Soit $g$ un automorphisme affin de $V$ tel que $gA=A'\in E$. Le segment $(gA, gB)=(A', B')$ de
$V$ possède la propriété 1\textsuperscript{o} de $(A, B)$. D'autre part, d'après
2\textsuperscript{o}, il existe un segment $(M, B')$ de $V$ tel que
\[ (A', B')\subset(M, B'); \qquad |p(M), p(A')|=r. \]
Posons
\[ p(M)=m; \qquad p(A')=a'; \qquad p(B')=b', \]
on a
\[ \frac{\overline{ma'}}{\overline{a'b'}}<k, \]
d'où
\[ \frac{r}{k}<|a'b'|. \]

4\textsuperscript{o} La démonstration va pouvoir s'achever. Soit $R$ le diamètre du compact
$p(E)\cup K$ de $\R^n$.

Choisissons $k$ tel que $\dfrac{r}{k}>R$.

Le segment $(A', B')$ construit dans 3\textsuperscript{o} vérifie la propriété
\[ A'\in E; \qquad p(B')\notin K. \]
D'où la proposition.
\end{proof}

\begin{prop}[2]
Soit $V$ une variété localement projective connexe balayable munie d'une projection $p$ sur
$S(n)$ compatible avec la structure projective [$V$ est un ouvert étalé sur $S(n)$]. Supposons de
plus qu'il existe un hémisphère ouvert $\mathring{H}$ de $S(n)$ tel que
\[ \overline{p(V)}\subset\mathring{H}. \]
Alors $p$ est une carte de $V$ sur l'intérieur d'un corps convexe de $S(n)$.
\end{prop}

\begin{proof}
L'enveloppe convexe de $p(V)$ [i.~e. le plus petit convexe de $\mathring{H}$ contenant $p(V)$] est
un convexe $W$, intérieur d'un corps convexe : nous allons montrer que $p$ est une carte de $V$
sur $W$.

% --- p. 258 ---
La démonstation se décompose en cinq points.
% [sic] « démonstation » : lire « démonstration ».

1\textsuperscript{o} Soit $g$ un automorphisme projectif de $V$, il existe un élément unique
$\gamma\in GS(n, \R)$ tel que
\[ p\circ g=\gamma\circ p. \]
$\gamma$ est un automorphisme projectif de $W$.

L'élément $\gamma$ a déjà été considéré \bzcf{en $b$ dans la démonstration du
théorème}{\ref{bzthm:c2s3t1_poche_balayable}}; que $\gamma$ soit un automorphisme de $W$ résulte
immédiatement de ce que $\gamma$ est un automorphisme de $p(V)$ dont $W$ est l'enveloppe convexe.

2\textsuperscript{o} Tout point $M\in V$ possède un voisinage $O_M$ tel que la restriction de $p$
à $O_M$ soit une carte sur une pseudo-boule (\bzcf{\textit{cf.} \S\,2, déf.
5}{\ref{bzdef:c2s2d5_nombre_recouvrement}}) de centre $p(M)$ et de rayon $r$ convenable. $O_M$ est
le lieu des points $N$ de $V$ tels qu'il existe un segment $(M, N)$ et que la pseudo-longueur du
segment $(p(M), p(N))$ dans $W$ soit inférieure à $r$. Par abus de langage, $O_M$ sera appelé une
pseudo-boule de centre $M$ et de rayon $r$.
% [sic] « § 2, déf. 5 » : lire « § 2, déf. 6 » (the pseudo-boule); the \ref names the printed
% déf. 5 (n(E, k)), per the convention of the header.

L'existence de $O_M$ résulte de ce que les pseudo-boules de centre $p(M)=m$ forment une base de
voisinage de $m$. Il est clair que tout point $N$ de $O_M$ possède la propriété énoncée;
réciproquement, si $n=p(N)$ appartient à la pseudo-boule de centre $m$ et que $N$ peut être joint
à $M$ par un segment de $V$, ce segment n'est autre que le relevé de $(m, n)$ dans $O_M$.

3\textsuperscript{o} Il existe un nombre $r$ positif, tel que tout point $M$ de $V$ possède un
voisinage $O_M$ pseudo-boule de centre $M$ et de rayon $r$.

Remarquons d'abord que, si $E$ est un compact de $V$, il existe $r(E)$ tel que tout point de $E$
soit centre d'une pseudo-boule de rayon $r(E)$ : cela résulte de 2\textsuperscript{o}, de la
continuité de la pseudo-longueur sur $W$, et de la compacité de $E$.

Soit $E$ un domaine générateur compact de $V$ : nous allons montrer que tout point de $V$ est
centre d'une pseudo-boule de rayon $r(E)$.

Soit $M\in V$ : il existe un automorphisme projectif $g$ de $V$ tel que $gM\in E$; soit $O'$ la
pseudo-boule de centre $gM$ et de rayon $r(E)$ : il résulte de 1\textsuperscript{o} que
$g^{-1}O'$ est une pseudo-boule $O$ de centre $M$ et de rayon $r(E)$.

4\textsuperscript{o} $p(V)=W$.

En effet, $p(V)$ est un ouvert de $W$ tel que, si $m\in p(V)$, la pseudo-boule de centre $m$ et de
rayon $r$ est incluse dans $p(V)$. Donc $p(V)=W$ (\bzcf{\textit{cf.} \S\,2, prop.
5}{\ref{bzprop:c2s2p5_ouvert_plein}}).

5\textsuperscript{o} $V$ est un revêtement de $W$.

Soit $m\in W$ : nous allons montrer que l'image réciproque par $p$ de la pseudo-boule $O$ de
centre $m$ et de rayon $r$ n'est autre que la réunion des pseudo-boules de rayon $r$ centrées aux
points de $p^{-1}(m)$, et que ces pseudo-boules sont deux à deux disjointes.

Que les pseudo-boules soient deux à deux disjointes résulte d'un résultat
% --- p. 259 ---
général de la théorie des ouverts étalés : deux sections d'un ouvert étalé (ici $V$) au-dessus
d'une même partie connexe de la base (ici la pseudo-boule de centre $m$ et de rayon $r$) sont
identiques ou d'intersection vide.

Soit $N\in p^{-1}(O)$ : nous allons montrer qu'il existe $M\in p^{-1}(m)$ tel que $N$ soit inclus
dans la pseudo-boule de centre $M$ et de rayon $r$.

En effet, soit $O'$ la pseudo-boule de centre $N$ et de rayon $r$ (une telle pseudo-boule existe
dans $V$ d'après 4\textsuperscript{o}) : on a $m\in p(O')$ : donc $O'$ contient un point $M$ de
$p^{-1}(m)$ : il est clair que $N$ est inclus dans la pseudo-boule de centre $M$ et de rayon $r$.
% [sic] « d'après 4° » : lire « d'après 3° » (point 3° gives the pseudo-boules of radius r).

Le point 5\textsuperscript{o} une fois établi, la proposition résulte trivialement de ce que $V$
est connexe, et $W$ simplement connexe.
\end{proof}

% (Chapter 3 « Compléments sur les convexes » begins on p. 259; not transcribed in this file.)

\end{document}
